% theme: evolution
\MajorQuestion{生物の進化}
\SetRowSpaceQanda

\Instruction{次の問いに答えなさい。}
\QQ{生物の形質が、長い年月をかけて世代を重ねるうちにしだいに変化することを\NamedBlank{evolution}という。}{進化}
\QQ{は虫類と鳥類の中間的な特徴をもつとされる化石生物を何というか。}{シソチョウ}
\QQ{魚類から進化したと考えられている脊椎動物は何類か。}{両生類}
\QQ{現在の形やはたらきは異なっていても、起源が同じであると考えられる器官を\NamedBlank{homologous_organ}という。}{相同器官}
\QQ{1859年に『種の起源』を発表した人物は誰か。}{ダーウィン}

\Instruction{\strut}
\QQ{植物は、水中で生活するものからどこで生活するものへ\RefBlank{evolution}してきたと考えられるか。}{陸上}
\QQ{魚類でありながら、肺をもち、ひれがあしのようになっている生物を何というか。}{ハイギョ}
\QQ{両生類から進化したと考えられている脊椎動物は何類か。}{は虫類}
\QQ{生物の特徴がさまざまであることを何というか。}{多様性}
\QQ{水中の光合成生物から進化したと考えられている植物を、2つ答えなさい。}{コケ植物、シダ植物}

\Instruction{\strut}
\QQ{脊椎動物は、水中からどこで生活できるように\RefBlank{evolution}してきたと考えられるか。}{陸上}
\QQ{哺乳類でありながら、くちばしをもち、卵を産む生物を何というか。}{カモノハシ}
\QQ{シダ植物から進化したと考えられている植物を何というか。}{裸子植物}
\QQ{両生類は、魚類とは虫類のどちらに近い特徴ももつか。}{両方}
\QQ{生物が生活する場所や、からだのつくりなどの共通する特徴を何というか。}{共通性}

\Instruction{\strut}
\QQ{裸子植物から進化したと考えられている植物を何というか。}{被子植物}
\QQ{イカの外とう膜の内側にある骨のようなもののように、はたらきを失って痕跡のみになっている器官を何というか。}{痕跡器官}
\QQ{は虫類から進化したと考えられている脊椎動物を、2つ答えなさい。}{鳥類、哺乳類}
\QQ{裸子植物と被子植物をまとめて何というか。}{種子植物}
\QQ{生物は、環境に適した形質をもつものへと長い年月をかけて変化してきた。この変化を何というか。}{進化}
